% A page with navigation landmarks, declared interaction states, and one
% backend-conditional sentence. The navs stand bare: each medium has its own
% answer — HTML keeps the <nav> landmarks, the PDF renders the unpinned menu
% as its document outline and drops the pinned control (viewport furniture),
% and the markdown twin drops both. One sentence is print-only, the honest
% ifbackend remainder. Invented content throughout.
\documentclass{article}
\pdfmeta{ title = "Field Notes", subject = "An invented page with a navigation landmark." }
\palette{ cobalt = #1D4ED8 }
\style{nav}{ hover = cobalt, focus = cobalt, motion = 150ms }
\begin{document}

\begin{nav}
\href{#first-light}{First Light} \href{#gathered-notes}{Gathered Notes} \href{#top}{Back to top}
\end{nav}

\section*{First Light}

An invented opening paragraph that appears on every surface.

\begin{ifbackend}{pdf}
This sentence is set only on the printed page.
\end{ifbackend}

\section*{Gathered Notes}

A second invented paragraph, also on every surface.

% The pinned return-to-top control: a labeled landmark (so two navs are
% distinguishable), fixed to the viewport corner, revealed by scroll where
% the platform can, and simply always visible with neither. Viewport
% furniture: the printed page has no viewport and ships nothing for it.
\begin{nav}[label = Return to top, pin = bottom right, offset = 1.5em, reveal = 300px]
\href{#top}{\faIcon[label = Return to top]{arrow-up}}
\end{nav}

\end{document}
